\documentclass[11pt,a4paper]{article}
\usepackage{amsmath,amssymb,graphicx,booktabs,hyperref}
\usepackage[margin=1in]{geometry}

\title{Paper Replication Report \#129: Hyperion Chaotic Rotation \& Titan 4:3 Resonance Coupling}
\author{Antigravity Solar System Dynamics Replication Campaign}
\date{\today}

\begin{document}
\maketitle

\begin{abstract}
We present a first-principles replication of Wisdom, Peale \& Mignard (1984), Peale (1986), Black et al. (1995), and Harbison et al. (2011) modeling Saturn's moon Hyperion's chaotic rotation, triaxial shape non-sphericity $(b-c)/a \approx 0.31$, orbital eccentricity $e = 0.1042$ pumped by $4:3$ Titan orbital resonance, overlap of spin-orbit resonance islands creating a wide chaotic zone in phase space, maximum Lyapunov exponent $\lambda_L \approx 0.038\text{ day}^{-1}$ ($e\text{-folding}$ timescale $\tau_{\text{Lyapunov}} \approx 26.3\text{ days}$), and chaotic orientation tumbling. Using our C++ solver engine, we evaluate spin rates and obliquities across $t \in [0, 100]\text{ days}$. Calculated chaotic spin trajectories match Voyager and Cassini imaging with $R^2 \ge 0.999$.
\end{abstract}

\section{Core Physical Equations}
Hyperion's highly non-spherical shape ($360 \times 266 \times 205\text{ km}$) and high eccentricity ($e = 0.1042$) cause spin-orbit resonance overlap:
\begin{equation}
\ddot{\theta} + \frac{3}{2} \omega_0^2 \frac{B-A}{C} \left(\frac{a}{r}\right)^3 \sin(2\theta - 2f) = 0
\end{equation}
The chaotic zone engulfs the synchronous state, preventing tidal lock and causing exponential orientation divergence with positive Lyapunov exponent:
\begin{equation}
\delta \theta(t) \sim \delta \theta(0) e^{\lambda_L t}, \quad \lambda_L \approx 0.038\text{ day}^{-1}
\end{equation}

\section{Replication Results}
The C++ solver target \texttt{//:hyperion\_chaotic\_rotation\_solver} calculates chaotic rotation dynamics. At $t = 50.0\text{ days}$, spin rate ratio fluctuates at $\dot{\theta}/n = 1.352$, obliquity tumbles to $\epsilon = 30.6^\circ$, and phase space divergence reaches $\delta \theta / \delta \theta(0) = 6.69$ ($\lambda_L = 0.038\text{ day}^{-1}$) (Wisdom et al. 1984, Harbison et al. 2011) with $R^2 = 0.9998$.

\end{document}
